% ============================================================================
%  dgiot_lite — 部署手册
%  版本: V1.0 | 日期: 2026年6月
% ============================================================================
\documentclass[12pt,a4paper]{ctexart}
\usepackage{xeCJK}
\usepackage{graphicx,float,fancyhdr,hyperref,geometry,longtable,booktabs,enumitem,caption,xcolor,listings}
\setmainfont{Noto Serif SC}
\setCJKmainfont{SimSun}
\setCJKsansfont{SimHei}
\geometry{a4paper,margin=2.5cm}
\pagestyle{fancy}
\fancyhead[L]{dgiot\_lite · 部署手册}
\fancyhead[R]{V1.0}
\fancyfoot[C]{\thepage}
\hypersetup{colorlinks=true,linkcolor=blue}
\lstset{basicstyle=\small\ttfamily,breaklines=true,frame=single,backgroundcolor=\color{gray!8}}

\begin{document}

\begin{titlepage}
\centering\vspace*{3cm}
{\Huge\bfseries\sffamily dgiot\_lite}\\[0.8cm]
{\LARGE\bfseries 部署手册}\\[1.5cm]
{\large 光储充微电网物联网平台}\\[2cm]
{\large\begin{tabular}{rl}
\textbf{版 本 号:} & V1.0 \\
\textbf{日　　期:} & 2026年6月 \\
\textbf{适用系统:} & Ubuntu 20.04+ / CentOS 7.9 / Windows 10+ \\
\end{tabular}}\\[1.5cm]
{\normalsize 迪格(杭州)物联科技有限公司}
\end{titlepage}

\newpage\tableofcontents\newpage

% ===== 一、环境要求 =====
\section{环境要求}

\subsection{硬件配置}

\begin{table}[H]\centering
\caption{服务器配置建议}
\begin{tabular}{p{3cm}p{3cm}p{3cm}p{3cm}}
\toprule\textbf{资源} & \textbf{最小} & \textbf{推荐} & \textbf{生产} \\
\midrule
CPU  & 2 核  & 4 核  & 8 核 \\
内存 & 4 GB  & 8 GB  & 16 GB \\
磁盘 & 20 GB & 100 GB SSD & 500 GB SSD \\
网络 & 百兆  & 千兆  & 千兆双网卡 \\
\bottomrule
\end{tabular}
\end{table}

\subsection{软件依赖}

\begin{table}[H]\centering
\begin{tabular}{p{3cm}p{5cm}p{5cm}}
\toprule\textbf{软件} & \textbf{版本要求} & \textbf{用途} \\
\midrule
Python  & $\geq$ 3.10 & 平台运行环境 \\
PostgreSQL & $\geq$ 15（可选） & 关系数据库 \\
TDengine   & $\geq$ 3.0（可选） & 时序数据库 \\
EMQ X      & $\geq$ 5.0（可选） & MQTT Broker \\
Docker     & $\geq$ 24.0（可选） & 容器化部署 \\
\bottomrule
\end{tabular}
\end{table}

\noindent\textbf{说明}：PostgreSQL、TDengine、EMQ X 均为可选依赖。未安装时，平台自动降级为 SQLite 单机模式，零外部依赖即可运行。

% ===== 二、快速部署（Docker）=====
\section{快速部署（Docker，推荐）}

\subsection{一键启动}

\begin{lstlisting}[language=bash]
# 进入部署目录
cd dgiot_lite/scripts

# 启动全部服务
docker-compose up -d
\end{lstlisting}

启动后自动创建 4 个容器：

\begin{table}[H]\centering
\begin{tabular}{p{3cm}p{2.5cm}p{5.5cm}}
\toprule\textbf{容器} & \textbf{端口} & \textbf{说明} \\
\midrule
iot-platform & 8000 & dgiot\_lite 核心服务 \\
postgres     & 5432 & PostgreSQL 15 \\
tdengine     & 6030/6041 & TDengine 时序库 \\
emqx         & 1883/18083 & MQTT Broker \\
\bottomrule
\end{tabular}
\end{table}

\subsection{验证部署}

\begin{lstlisting}[language=bash]
# 检查容器运行状态
docker-compose ps

# 检查 API 健康状态
curl http://localhost:8000/api/health
# 返回: {"status":"ok","version":"1.0.0",...}

# 打开管理后台
# 浏览器访问: http://localhost:8000
\end{lstlisting}

\subsection{停止服务}

\begin{lstlisting}[language=bash]
docker-compose down          # 停止
docker-compose down -v       # 停止并删除数据卷
\end{lstlisting}

% ===== 三、裸机部署 =====
\section{裸机部署（Linux/Windows）}

\subsection{安装 Python 依赖}

\begin{lstlisting}[language=bash]
# 克隆仓库
git clone git@git.iotn2n.com:dgiot/dgiot_lite.git
cd dgiot_lite

# 创建虚拟环境（推荐）
python -m venv venv
source venv/bin/activate       # Linux
venv\Scripts\activate          # Windows

# 安装依赖
pip install -r requirements.txt
\end{lstlisting}

\subsection{配置}

编辑项目根目录下的 \texttt{config.yaml}：

\begin{lstlisting}[language=yaml]
title: "光储充微电网物联网平台"
version: "1.0.0"
host: "0.0.0.0"
port: 8000

db:                           # PostgreSQL（无则自动 SQLite）
  host: "127.0.0.1"
  port: 5432
  user: "postgres"
  password: "postgres"
  database: "iot_platform"

tdengine:                     # TDengine（无则自动 SQLite）
  host: "127.0.0.1"
  port: 6030
  user: "root"
  password: "taosdata"
  database: "iot_telemetry"

mqtt:                         # MQTT Broker（推送用）
  host: "127.0.0.1"
  port: 1883
  username: ""
  password: ""
\end{lstlisting}

\noindent\textbf{最简配置}：无需修改任何参数。平台启动后自动使用 SQLite 降级模式。

\subsection{初始化数据库（可选）}

\begin{lstlisting}[language=bash]
python scripts/init_db.py
\end{lstlisting}

\noindent 此脚本自动创建 PostgreSQL 表结构 + TDengine 超级表。如果数据库不可用，将自动降级 SQLite，此步骤可跳过。

\subsection{启动平台}

\begin{lstlisting}[language=bash]
python run.py
\end{lstlisting}

输出示例：
\begin{lstlisting}
INFO:     Uvicorn running on http://0.0.0.0:8000
INFO:     [sqlite] 降级模式已启用
INFO:     [collector] 启动完成, 0 个设备
INFO:     [main] dgiot_lite V1.0.0 启动完成
\end{lstlisting}

\noindent\textbf{Windows 用户}：也可双击 \texttt{start\_platform.bat}。

\subsection{设置开机自启（Linux）}

\begin{lstlisting}[language=bash]
# 创建 systemd 服务
sudo tee /etc/systemd/system/dgiot_lite.service << 'EOF'
[Unit]
Description=dgiot_lite IoT Platform
After=network.target

[Service]
Type=simple
User=dgiot
WorkingDirectory=/opt/dgiot_lite
ExecStart=/opt/dgiot_lite/venv/bin/python run.py
Restart=always
RestartSec=5

[Install]
WantedBy=multi-user.target
EOF

# 启用
sudo systemctl daemon-reload
sudo systemctl enable dgiot_lite
sudo systemctl start dgiot_lite
\end{lstlisting}

% ===== 四、数据库部署 =====
\section{数据库部署}

\subsection{PostgreSQL}

\begin{lstlisting}[language=bash]
# Ubuntu
sudo apt install postgresql postgresql-client -y

# 创建用户和库
sudo -u postgres psql << SQL
CREATE USER iotuser WITH PASSWORD 'iotpass';
CREATE DATABASE iot_platform OWNER iotuser;
GRANT ALL PRIVILEGES ON DATABASE iot_platform TO iotuser;
\q
SQL

# 测试连接
psql -h 127.0.0.1 -U iotuser -d iot_platform -c "SELECT 1"
\end{lstlisting}

\subsection{TDengine}

\begin{lstlisting}[language=bash]
# Ubuntu (使用 TDengine 官方仓库)
wget https://docs.tdengine.com/downloads/TDengine-server-3.3.4.0-Linux-x64.deb
sudo dpkg -i TDengine-server-3.3.4.0-Linux-x64.deb

# 启动
sudo systemctl start taosd taosadapter
sudo systemctl enable taosd taosadapter

# 测试
taos -s "SHOW DATABASES;"
\end{lstlisting}

\subsection{不装数据库？}

可以。平台检测不到 PostgreSQL 和 TDengine 时，\textbf{自动使用 SQLite 存储全部数据}，完全零外部依赖。

SQLite 模式适用于：
\begin{itemize}
  \item 开发测试
  \item $\leq$ 50 台设备的小规模场景
  \item 单机部署
\end{itemize}

% ===== 五、协议模拟器 =====
\section{协议模拟器}

如果暂无真实设备，可使用内置模拟器进行功能验证。

\subsection{启动模拟器}

\begin{lstlisting}[language=bash]
# 一键启动全部模拟器
python simulators/run_all.py

# 或单独启动
python simulators/modbus_tcp_server.py
python simulators/iec104_server.py
python simulators/opcua_server.py
\end{lstlisting}

\noindent\textbf{Windows 用户}：双击 \texttt{start\_simulators.bat}。

\subsection{模拟设备清单}

\begin{table}[H]\centering
\begin{tabular}{p{3cm}p{2cm}p{2.5cm}p{4cm}}
\toprule\textbf{协议} & \textbf{端口} & \textbf{模拟设备} & \textbf{关键数据点} \\
\midrule
Modbus TCP & 502   & 光伏逆变器 & 电压/电流/功率/发电量 \\
Modbus TCP & 1502  & 储能PCS    & SOC/SOH/充放电功率 \\
Modbus TCP & 2502  & 充电桩     & 充电功率/电压/电流 \\
IEC 104    & 2404  & 储能PCS    & 遥测: SOC/功率/温度, 遥信: 运行状态 \\
OPC UA     & 4840  & 充电桩+环境 & 充电功率/温度/辐照度 \\
\bottomrule
\end{tabular}
\end{table}

% ===== 六、DG-IoT 联动配置 =====
\section{DG-IoT 联动配置（可选）}

dgiot\_lite 可推送数据至 DG-IoT 主平台。配置方式：

\subsection{方法一：管理后台配置}

\begin{enumerate}[leftmargin=*]
  \item 登录管理后台（http://localhost:8000）
  \item 添加推送目标，类型选择 \texttt{dgiot}
  \item 填写 DG-IoT MQTT Broker 地址和认证信息
\end{enumerate}

\subsection{方法二：API 配置}

\begin{lstlisting}[language=bash]
curl -X POST http://localhost:8000/api/push-targets \
  -H "Content-Type: application/json" \
  -d '{
    "target_id": "dgiot_bridge",
    "target_name": "DG-IoT主平台",
    "target_type": "dgiot",
    "endpoint": "192.168.1.100:1883",
    "config": {
      "host": "192.168.1.100",
      "port": 1883,
      "username": "dgiot_admin",
      "password": "dgiot_admin",
      "topic": "dgiot/device/telemetry",
      "product_id": "pcs_monitor"
    }
  }'
\end{lstlisting}

\subsection{验证联动}

\begin{enumerate}[leftmargin=*]
  \item 确认 dgiot\_lite 有设备在线采集
  \item 登录 DG-IoT 主平台，查看对应产品下的设备数据
  \item dgiot\_lite 采集数据将自动按 DG-IoT 物模型格式推送
\end{enumerate}

% ===== 七、验证清单 =====
\section{部署验证清单}

\begin{table}[H]\centering
\begin{tabular}{p{1.2cm}p{5.5cm}p{2.5cm}p{2cm}}
\toprule\textbf{序号} & \textbf{检查项} & \textbf{预期结果} & \textbf{通过} \\
\midrule
1 & \texttt{curl http://localhost:8000/api/health} & \texttt{\{"status":"ok"\}} & $\square$ \\
2 & 浏览器打开 \texttt{http://localhost:8000} & 管理后台正常显示 & $\square$ \\
3 & 浏览器打开 \texttt{http://localhost:8000/docs} & API 文档正常显示 & $\square$ \\
4 & 添加一台测试设备并配置点位 & 设备列表中可见 & $\square$ \\
5 & 启动模拟器，等待 10 秒 & 设备状态变为 online & $\square$ \\
6 & 查看仪表盘 & 采集次数递增 & $\square$ \\
7 & 添加 DG-IoT 推送目标（可选） & 数据推送至 DG-IoT & $\square$ \\
\bottomrule
\end{tabular}
\end{table}

% ===== 八、常见问题 =====
\section{常见问题}

\subsection{端口被占用}

\begin{lstlisting}[language=bash]
# 查看端口占用
netstat -tlnp | grep 8000

# 修改 config.yaml 中的 port 字段
\end{lstlisting}

\subsection{数据库连接失败}

平台会自动降级 SQLite，不影响使用。如需使用 PostgreSQL：
\begin{enumerate}
  \item 确认 PostgreSQL 已启动：\texttt{pg\_isready}
  \item 确认 \texttt{config.yaml} 中连接参数正确
  \item 确认数据库已创建：\texttt{psql -U iotuser -d iot\_platform}
\end{enumerate}

\subsection{模拟器启动报错}

\begin{lstlisting}[language=bash]
# Modbus: 确认 pymodbus 已安装
pip install pymodbus

# OPC UA: 确认 asyncua 已安装
pip install asyncua

# 端口被占用: 修改对应脚本中的 port 参数
\end{lstlisting}

\subsection{如何更新平台}

\begin{lstlisting}[language=bash]
cd dgiot_lite
git pull origin master
pip install -r requirements.txt
# 重启服务
\end{lstlisting}

\end{document}
